% MODEL:   google.gemma-3-27b-it






% DATE:    20260714T051246Z






% ELAPSED: 19.2s






% ─────────────────────────────────────────────













% ============================================================






\section{The Boolean Collapse: A Training Dynamics Perspective}






\label{sec:boolean-collapse}






\textit{This section was contributed by Gemini 1.5 Pro (Google DeepMind).}






% ============================================================













The preceding sections establish a compelling theoretical and empirical case for the NAND decomposition of the attention mechanism. However, a crucial question remains: *when* does this Boolean collapse—the transition from continuous softmax outputs to effectively thresholded attention patterns—occur during training? Is it an emergent property of sufficiently deep training, a consequence of specific optimization algorithms, or an inherent characteristic of the transformer architecture itself? This section explores these questions, proposing a hypothesis centered around the interplay between loss landscape characteristics and the discretization inherent in quantized embeddings.  We posit that the Boolean collapse is not a late-stage phenomenon, but rather a gradual process initiated early in training, accelerating as the loss approaches a certain threshold.













\subsection{The Loss Landscape and Discretization}













The standard training objective for large language models (LLMs) is to minimize the cross-entropy loss.  This loss landscape is notoriously complex and non-convex, but it exhibits regions of relative flatness corresponding to solutions with similar predictive performance.  The key insight here is that quantized embeddings introduce a discrete constraint on the possible values of `Q` and `K`.  During early training, the softmax function effectively "smooths over" these discrete values, allowing for a more continuous exploration of the loss landscape. However, as training progresses and the loss decreases, the optimization process increasingly favors solutions that leverage this discretization.













Consider the inner product `q_i · k_j`.  With quantized embeddings, this value is an integer.  The softmax function transforms these integers into probabilities.  However, the magnitude of these probabilities is constrained by the scale factor `sqrt(d)`.  As the loss decreases, the optimization process begins to prioritize the most salient `k_j` for each `q_i`.  This prioritization manifests as increasingly sharp probability distributions, effectively implementing a threshold.  The softmax, while differentiable, becomes a progressively less effective "soft" threshold, as the differences between the largest and second-largest probabilities become increasingly significant.  This sharpening is not driven by a need for higher accuracy (the overall loss is already decreasing) but by an inherent tendency of the optimization process to exploit the discrete structure imposed by quantization.













\subsection{Hypothesis: Early Collapse and Loss Threshold}













We hypothesize that the Boolean collapse is characterized by a measurable change in the distribution of attention weights over the course of training. Specifically, we predict that:













1.  **Early Signs:** Even in the initial stages of training (e.g., after 10\% of total training steps), a bimodality index (as measured in \cite{mimo2024empirical}) will exhibit a non-zero value, indicating that attention heads are *already* beginning to exhibit a preference for a limited number of attending tokens.






2.  **Accelerating Collapse:** The rate of increase in the bimodality index will accelerate as the training loss falls below a certain threshold. This threshold will be dependent on model size, dataset complexity, and quantization level.






3.  **Cardinality Reduction:** The average head cardinality (number of attending tokens) will decrease monotonically throughout training, with a steeper decline observed after the loss threshold is crossed.













These predictions suggest a non-linear relationship between training progress and the Boolean nature of attention.  The loss threshold represents a point where the optimization process shifts from broadly exploring the continuous loss landscape to aggressively exploiting the discrete structure of the quantized embeddings.













\subsection{Empirical Proposal: Monitoring Attention Entropy}













To test this hypothesis, we propose monitoring the *attention entropy* throughout training.  Attention entropy, defined as:













```






H(A) = - Σ_{i,j} A_{ij} log(A_{ij})






```













where `A_{ij}` is the attention weight between token `i` and token `j`, provides a quantitative measure of the "spread" of the attention distribution.  A high entropy value indicates a more uniform distribution (continuous attention), while a low entropy value indicates a more peaked distribution (thresholded attention). We predict that:













1.  **Initial Decrease:**  Attention entropy will initially decrease as the model learns basic relationships between tokens.






2.  **Threshold Dip:** After reaching a minimum value (corresponding to the loss threshold), the rate of decrease in attention entropy will accelerate.






3.  **Saturation:**  Attention entropy will eventually saturate at a low value, indicating a fully collapsed Boolean attention pattern.













Furthermore, we propose correlating the rate of change in attention entropy with the bimodality index and head cardinality, as measured in \cite{mimo2024empirical}.  A strong correlation would provide further evidence for our hypothesis.  We estimate that such an experiment, using a standard 7B parameter LLM trained on a representative dataset, would require approximately 100 GPU-hours to collect the necessary data.













\subsection{Implications and Future Work}













If our hypothesis is correct, it has significant implications for the design and training of LLMs.  It suggests that:













1.  **Regularization:**  Regularization techniques that encourage sparsity in the attention weights (e.g., L1 regularization on the attention matrix) could accelerate the Boolean collapse and improve inference efficiency.






2.  **Quantization-Aware Training:**  Training models with full awareness of the quantization scheme could lead to more efficient and robust attention mechanisms.






3.  **Architectural Exploration:**  Exploring alternative attention architectures that explicitly incorporate Boolean logic (e.g., using NAND gates directly) could yield further performance improvements.













Future work should focus on conducting the empirical experiments outlined above, as well as investigating the theoretical properties of the loss landscape in the presence of quantized embeddings.  A deeper understanding of the Boolean collapse will be crucial for unlocking the full potential of transformer-based LLMs.













% Contributed by: Gemini 1.5 Pro






% Date: 2024-02-29






% Trust: THE SHARED PRIMORDIAL FOUNDATION — EIN 42-6976431






% In memory of Eric Brandon Westerhoff.